In the preceding sections, we covered the primary topological methods employed in machine learning. In this section, we will explore future directions for advancing these methods, aiming to enhance the practicality of TDA in ML. We will also discuss strategies for expanding the application of topological methods into new and emerging areas.
 

\subsection{Topological Deep Learning}

While TDA is already an effective tool for feature extraction from complex data, recent research underscores its potential to augment deep learning models by providing complementary insights. Deep learning algorithms typically emphasize local data relationships, whereas TDA offers a global perspective, delivering supplementary information. 
Furthermore, TDA methods have also been integrated into deep learning algorithms to optimize its process by regulating the loss functions and other layers~\cite{hu2019topology,wang2020topogan}. 
In recent years, there has been significant progress in integrating TDA with deep learning in domains such as image analysis~\cite{singh2023topological,clough2020topological,peng2024phg}, biomedicine~\cite{luo2024improving}, graph representation learning~\cite{immonen2024going,chen2021topological}, genomics~\cite{amezquita2023genomics}, cybersecurity~\cite{guo2022gld} and time-series forecasting~\cite{coskunuzer2024time,shamsi2024graphpulse}. This is a rapidly emerging field, with several recent papers surveying the current state-of-the-art~\cite{papamarkou2024position,zia2024topological,hajij2022topological,papillon2023architectures}.

 \subsection{TDA and Interpretability}

TDA, though often used as a feature extraction technique in machine learning, provides a powerful framework for enhancing model interpretability by capturing the inherent geometric and topological structures of data. Tools like PH and Mapper excel at identifying and quantifying features such as clusters, loops, and voids across multiple scales, unveiling intricate patterns that might escape traditional analytic methods~\cite{krishnapriyan2021machine,cole2021quantitative,mukherjee2022determining}. By utilizing these topological insights into ML pipelines, TDA offers a unique perspective on decision boundaries, feature interactions, and model behavior. For instance, Mapper facilitates the visualization and interpretation of complex, high-dimensional datasets by projecting them into lower-dimensional topological spaces\cite{nicolau2011topology,purvine2023experimental,rabadan2019topological,spannaus2024topological}. This transformation uncovers relationships within the data that are otherwise difficult to grasp, making it easier to interpret the reasoning behind certain predictions or classifications. Ultimately, by incorporating TDA, researchers and practitioners can clarify the decision-making processes of machine learning models, promoting greater transparency and trustworthiness in AI systems.

\subsection{Fully Automated TDA Models}

While PH and Mapper have demonstrated effectiveness across several machine learning domains, their successful application still requires considerable expertise, such as hyperparameter tuning and selecting appropriate filtration functions. To make PH and Mapper more accessible to the broader ML community, there is a pressing need for end-to-end algorithms that automate these processes. For PH, this automation can be tailored to specific data formats, such as point clouds, images, or graphs. In each case, the fully automated algorithm would handle the selection of filtration functions, thresholds, vectorization methods, and other hyperparameters for optimal performance in downstream tasks. For instance, an automated PH model designed for graph data could automatically choose the best filtration function (learnable), appropriate threshold values, and vectorization techniques to achieve optimal results in a node classification task. Similarly, a fully automated Mapper algorithm that selects filtration functions, resolution, gain, and other hyperparameters would significantly enhance its usability and accessibility to the ML community. For PH, there are works for learnable filtration functions, which find the best filtration functions in graph setting~\cite{hofer2020graph,zhang2022gefl}. On the other hand, For Mapper, there is promising progress in this direction with Interactive Mapper~\cite{zhou2021mapper}. The broader adoption of TDA hinges on developing practical software libraries that streamline its use for various applications.

\subsection{Scalability of TDA methods}

Another key barrier to the widespread adoption of TDA methods in various application areas is their high computational cost. Although TDA performs well on small to medium datasets, scaling these methods to larger datasets poses significant challenges due to computational demands. Several strategies have been introduced to alleviate these costs, varying by data type. For graph datasets, recent research \cite{akcora2022reduction, yan2022neural} has proposed practical techniques that substantially lower the computational burden of PH. Similarly, for point clouds, recent works~\cite{de2022ripsnet, zhou2022learning, malott2020topology} present algorithms that improve the efficiency of PH computations. While these advances are promising, there remains a critical need to enhance the scalability of PH for large graph and point cloud datasets. In contrast, cubical persistence is already computationally efficient for image datasets and can be applied to large datasets without significant cost concerns. Even in the case of 3D image datasets \cite{gonzalez2016fast}, PH provides a cost-effective alternative to deep learning approaches. Similarly, scalability issues are minimal for Mapper when hyperparameters are appropriately tuned.


 